% ============================================================================
% Fichier  : partie1/P01_C02_demarche_enquete.tex : FORN (LPI342)
% Sources : ISO/IEC 27043 ; NIST SP 800-86 ; ACPO. Opération TSANAGA.
% ============================================================================

\chapter{La Démarche d'Enquête}
\label{chap:demarche}
\chaptermark{La démarche d'enquête}

\epigraph{%
    « Une enquête n'est pas une suite d'outils~;\\
    c'est une suite de décisions documentées. »%
}{Principe du cours}

\niveauChapitre{Fondamental}{Chapitre~\ref{chap:forensique}.}

\begin{boxobjectifs}
\begin{itemize}
    \item Décrire les phases d'une enquête numérique.
    \item Comprendre la planification et l'autorisation.
    \item Situer identification, acquisition, analyse, restitution.
    \item Distinguer constatation et interprétation.
    \item Documenter chaque phase.
\end{itemize}
\end{boxobjectifs}

\bigskip

% ============================================================================
\section*{Question directrice : par où commence une enquête sur les supports de M. NDONGO ?}
% ============================================================================

Devant les scellés, la tentation est d'« aller voir ». L'enquête impose un
ordre~: planifier, identifier, acquérir, analyser, restituer. Ce chapitre
pose cette démarche.

\medskip
\textbf{Plan du chapitre.}
\begin{enumerate}
    \item Les phases (section~\ref{sec:c02-phases}).
    \item Planifier et autoriser (section~\ref{sec:c02-planifier}).
    \item Constatation et interprétation (section~\ref{sec:c02-constat}).
    \item Documenter (section~\ref{sec:c02-documenter}).
\end{enumerate}

\separateur

% ============================================================================
\section{Les Phases}
\label{sec:c02-phases}
% ============================================================================

\begin{tableaucours}{L{3.4cm}L{8.8cm}}{%
    \tableauHeader{Phase} &
    \tableauHeader{Ce qu'on y fait}
}
Planifier & Cadre, autorisation, périmètre, outils \\
Identifier & Repérer et préserver les sources de preuve \\
Acquérir & Dupliquer, calculer les empreintes, tracer \\
Analyser & Examiner la copie, interpréter, recouper \\
Restituer & Rapport, présentation, conservation \\
\end{tableaucours}
\vspace{-0.8em}
\begin{center}{\small\itshape Tableau~2.1. Les phases d'une enquête~\cite{iso27043,nist80086}}\end{center}

% ============================================================================
\section{Planifier et Autoriser}
\label{sec:c02-planifier}
% ============================================================================

\begin{boxjuris}
Une enquête forensique suppose un fondement~: réquisition judiciaire, mandat,
ou autorisation écrite du responsable pour une enquête interne. Au Cameroun,
accéder à un support sans titre expose aux sanctions de la
\loicam{loi n°~2010/012} (Chapitre~\ref{chap:droit}). On vérifie son titre
avant d'agir.
\end{boxjuris}

% ============================================================================
\section{Constatation et Interprétation}
\label{sec:c02-constat}
% ============================================================================

\begin{boxalert}
\textbf{Violation typique.} Écrire « M.~NDONGO a volé les données » dès la
première trace.

\smallskip
\textbf{Règle~:} on constate (« le fichier clients a été copié sur une clé le
3~janvier »), puis on interprète séparément, avec un niveau de confiance. Le
rapport sépare toujours le fait établi de l'hypothèse.
\end{boxalert}

% ============================================================================
\section{Documenter}
\label{sec:c02-documenter}
% ============================================================================

\begin{boxscenario}{Opération TSANAGA. Le journal d'enquête}
\textbf{Constatation.} Dès la réception des scellés, Mme~NKENG ouvre un
journal d'enquête~: chaque acte y est horodaté, décrit, signé.

\medskip
\textbf{Analyse} (confiance élevée). Le journal n'est pas une formalité~:
c'est lui qui rendra l'enquête reproductible et opposable. Une analyse non
documentée, même exacte, sera contestée au tribunal.
\end{boxscenario}

\begin{boxmetier}{Enquêteur en cybercriminalité}
\textbf{Suivre la démarche sous pression.} Même pressé par le parquet,
l'enquêteur ne saute pas de phase~: planifier et identifier avant d'acquérir,
acquérir avant d'analyser. L'ordre protège la preuve.
\end{boxmetier}

\separateur

\begin{boxresume}
\textbf{Ce qu'il faut retenir de ce chapitre~:}
\begin{itemize}
    \item \textbf{Cinq phases}~: planifier, identifier, acquérir, analyser,
          restituer.
    \item \textbf{On vérifie son titre} avant d'agir.
    \item \textbf{On constate, puis on interprète}, séparément.
    \item \textbf{On documente} chaque acte dans un journal d'enquête.
    \item \textbf{L'ordre des phases} protège la preuve.
\end{itemize}
\end{boxresume}

\begin{boxinfo}
\textbf{Points de contrôle de l'élève.} Je sais~: énoncer les phases~;
expliquer le fondement d'une enquête~; séparer constatation et
interprétation~; tenir un journal d'enquête.
\end{boxinfo}

\bigskip

\begin{boxexo}[title={\ding{45}\quad Exercice~2.1. Les phases \niveaubase}]
Citez les cinq phases d'une enquête numérique dans l'ordre.
\end{boxexo}

\begin{boxexo}[title={\ding{45}\quad Exercice~2.2. Le titre \niveaubase}]
Quel fondement autorise l'enquête sur les supports de M.~NDONGO~?
\end{boxexo}

\begin{boxexo}[title={\ding{45}\quad Exercice~2.3. Constat ou interprétation \niveaubase}]
Classez~: « le fichier a été copié le 3~janvier »~; « M.~NDONGO est
coupable ».
\end{boxexo}

\begin{boxexo}[title={\ding{45}\quad Exercice~2.4. Sauter une phase \niveaumoyen}]
Que risque-t-on à analyser avant d'acquérir~?
\end{boxexo}

\begin{boxexo}[title={\ding{45}\quad Exercice~2.5. Journal \niveaumoyen}]
Que consigne-t-on dans un journal d'enquête~?
\end{boxexo}

\begin{boxexo}[title={\ding{45}\quad Exercice~2.6. Confiance \niveaumoyen}]
Pourquoi assortir une interprétation d'un niveau de confiance~?
\end{boxexo}

\begin{boxexo}[title={\ding{45}\quad Exercice~2.7. Pression \niveaumoyen}]
Le parquet veut un résultat vite. Comment tenir la démarche~?
\end{boxexo}

\begin{boxexo}[title={\ding{45}\quad Exercice~2.8. Reproductible \niveaumoyen}]
Pourquoi une enquête doit-elle être reproductible~?
\end{boxexo}

\begin{boxexo}[title={\ding{45}\quad Exercice~2.9. Plan d'enquête \niveauexpert}]
Rédigez le plan d'enquête de l'opération TSANAGA, phase par phase.
\end{boxexo}

\begin{boxexo}[title={\ding{45}\quad Exercice~2.10. Modèle de journal \niveauexpert}]
Concevez le modèle de journal d'enquête du cabinet NYONG SÉCURITÉ.
\end{boxexo}

\bigskip

\begin{boxinfo}
\textbf{Pour aller plus loin.}
\begin{itemize}
    \item ISO/IEC 27043~\cite{iso27043}~; NIST SP~800-86~\cite{nist80086}~;
          ACPO~\cite{acpo2012}.
\end{itemize}
\end{boxinfo}

\begin{boxinfo}
\textbf{Transition.} La démarche n'a de valeur que portée par une éthique~:
le Chapitre~\ref{chap:ethique} traite de la neutralité et de l'impartialité.
\end{boxinfo}
